Bounce-averaged spatial flux divergence

In this section, calculations are restricted to the bounce averaged
anomalous radial transport equation that is added to the Fokker-Planck
equation in an heuristic manner. Here, momentum dynamics is omitted for sake
of clarity. It is considered that the anomalous transport equation may be
cast in a conservative form as for the Fokker-Planck equation. Therefore,
the problem is reduced to the calculation of the bounce averaged divergence
of the spatial flux $\mathbf{S}_{\mathbf{r}}$ 
\begin{equation}
\left\{ \nabla _{\mathbf{r}}\cdot \mathbf{S}_{\mathbf{r}}\right\} =\left\{ 
\frac{1}{J_{\mathbf{r}}}\frac{\partial }{\partial u^{i}}\left( J_{\mathbf{r}
}S_{\mathbf{r}}^{i}\right) \right\}
\end{equation}

For this purpose, the contravariant vector basis $\left( \widehat{r},%
\widehat{\theta },\widehat{\phi }\right) $ must be used, since it is
associated to length coordinates in configuration. Indeed, $du^{i}$ are
elementary lengths along the coordinate curves. Using general identities 
\begin{equation}
\left\{ 
\begin{array}{l}
\partial u^{r}=\left| \nabla r\right| ^{-1}\partial r \\ 
\partial u^{\theta }=\left| \nabla \theta \right| ^{-1}\partial \theta \\ 
\partial u^{\phi }=\left| \nabla \phi \right| ^{-1}\partial \phi
\end{array}
\right.
\end{equation}
with 
\begin{equation}
\left\{ 
\begin{array}{l}
\left| \nabla r\right| ^{-1}=1 \\ 
\left| \nabla \theta \right| ^{-1}=r \\ 
\left| \nabla \phi \right| ^{-1}=R
\end{array}
\right.
\end{equation}
and the Jacobian $J_{\mathbf{r}}=\left| \nabla r\times \nabla \theta \cdot
\nabla \phi \right| ^{-1}=rR,$ it comes 
\begin{equation}
\frac{\partial }{\partial u^{r}}\left( J_{\mathbf{r}}S_{\mathbf{r}
}^{r}\right) =\frac{\partial }{\partial r}\left( rRS_{\mathbf{r}}^{r}\right)
\end{equation}
\begin{equation}
\frac{\partial }{\partial u^{\theta }}\left( J_{\mathbf{r}}S_{\mathbf{r}%
}^{\theta }\right) =\frac{1}{r}\frac{\partial }{\partial \theta }\left( rRS_{%
\mathbf{r}}^{\theta }\right)
\end{equation}
\begin{equation}
\frac{\partial }{\partial u^{\phi }}\left( J_{\mathbf{r}}S_{\mathbf{r}%
}^{\phi }\right) =\frac{1}{R}\frac{\partial }{\partial \phi }\left( rRS_{%
\mathbf{r}}^{\phi }\right) =0
\end{equation}
taking benefits of the axisymmetry. Finally, 
\begin{eqnarray}
\nabla _{\mathbf{r}}\cdot \mathbf{S}_{\mathbf{r}} &=&\frac{1}{rR}\left[ 
\frac{\partial }{\partial r}\left( rRS_{\mathbf{r}}^{r}\right) +\frac{1}{r}%
\frac{\partial }{\partial \theta }\left( rRS_{\mathbf{r}}^{\theta }\right)
\right]  \nonumber \\
&=&\frac{1}{rR}\frac{\partial }{\partial r}\left( rRS_{\mathbf{r}
}^{r}\right) +\frac{1}{r^{2}R}\frac{\partial }{\partial \theta }\left( rRS_{%
\mathbf{r}}^{\theta }\right)
\end{eqnarray}

In order to perform bounce-averaging calculations, it is more appropriate to
express the components of $\mathbf{S}_{\mathbf{r}}$ in the covariant vector
basis $\left( \widehat{\psi },\widehat{s},\widehat{\phi }\right) $ since in
the thin banana limit, particle trajectories remain confined in a magnetic
surface labelled by the poloidal magnetic flux $\psi $ . Therefore, 
\begin{equation}
\mathbf{S}_{\mathbf{r}}=S_{\mathbf{r}}^{r}\widehat{r}+S_{\mathbf{r}}^{\theta
}\widehat{\theta }=S_{\mathbf{r}}^{\psi }\widehat{\psi }+S_{\mathbf{r}}^{s}%
\widehat{s}
\end{equation}
where the covariant vector basis may be deduced from the contravariant one
by the rotation transform characterized by the matrix $R$

\begin{equation}
\left( 
\begin{array}{l}
\widehat{\psi } \\ 
\widehat{s}
\end{array}
\right) =R\left( 
\begin{array}{l}
\widehat{r} \\ 
\widehat{\theta }
\end{array}
\right)
\end{equation}
according to the relation

\begin{equation}
R\left( 
\begin{array}{ll}
\widehat{\psi }\widehat{\cdot r} & \widehat{\psi }\cdot \widehat{\theta } \\ 
\widehat{s}\widehat{\cdot r} & \widehat{s}\cdot \widehat{\theta }
\end{array}
\right) =\left( 
\begin{array}{ll}
\cos \alpha & -\sin \alpha \\ 
\sin \alpha & \cos \alpha
\end{array}
\right)
\end{equation}

Flux vector coordinates relations are then 
\begin{equation}
\left( 
\begin{array}{l}
S_{\mathbf{r}}^{r} \\ 
S_{\mathbf{r}}^{\theta }
\end{array}
\right) =R^{-1}\left( 
\begin{array}{l}
S_{\mathbf{r}}^{\psi } \\ 
S_{\mathbf{r}}^{s}
\end{array}
\right)
\end{equation}
and,

\begin{equation}
\nabla _{\mathbf{r}}\cdot \mathbf{S}_{\mathbf{r}}=\frac{1}{rR}\frac{\partial 
}{\partial r}\left( rR\left[ S_{\mathbf{r}}^{\psi }\cos \alpha +S_{\mathbf{r}%
}^{s}\sin \alpha \right] \right) +\frac{1}{r^{2}R}\frac{\partial }{\partial
\theta }\left( rR\left[ -S_{\mathbf{r}}^{\psi }\sin \alpha +S_{\mathbf{r}
}^{s}\cos \alpha \right] \right)
\end{equation}

From the relations between the coordinate systems $\left( r,\theta \right) $
and $\left( \psi ,s\right) ,$ 
\begin{equation}
\left( 
\begin{array}{l}
\partial /\partial r \\ 
\partial /\partial \theta
\end{array}
\right) =\left( 
\begin{array}{ll}
\partial \psi /\partial r & \partial s/\partial r \\ 
\partial \psi /\partial \theta & \partial s/\partial \theta
\end{array}
\right) \cdot \left( 
\begin{array}{l}
\partial /\partial \psi \\ 
\partial /\partial s
\end{array}
\right)
\end{equation}
with 
\begin{eqnarray}
\left. \frac{\partial \psi }{\partial r}\right| _{s} &=&\frac{\nabla \psi
\cdot \widehat{\psi }}{\nabla r\cdot \widehat{r}}\left( \widehat{\psi }\cdot 
\widehat{r}\right) =\left( \nabla \psi \cdot \widehat{\psi }\right) \cos
\alpha  \nonumber \\
\left. \frac{\partial \psi }{\partial \theta }\right| _{s} &=&\frac{\nabla
\psi \cdot \widehat{\psi }}{\nabla \theta \cdot \widehat{\theta }}\left( 
\widehat{\psi }\cdot \widehat{\theta }\right) =-\left( \nabla \psi \cdot 
\widehat{\psi }\right) r\sin \alpha  \nonumber \\
\left. \frac{\partial s}{\partial r}\right| _{\psi } &=&\frac{\nabla s\cdot 
\widehat{s}}{\nabla r\cdot \widehat{r}}\left( \widehat{s}\cdot \widehat{r}%
\right) =\sin \alpha  \nonumber \\
\left. \frac{\partial s}{\partial \theta }\right| _{\psi } &=&\frac{\nabla
s\cdot \widehat{s}}{\nabla \theta \cdot \widehat{\theta }}\left( \widehat{s}%
\cdot \widehat{\theta }\right) =r\cos \alpha
\end{eqnarray}
the transport equation may be recast in a more natural form for the
bounce-averaging procedure 
\begin{eqnarray}
\nabla _{\mathbf{r}}\cdot \mathbf{S}_{\mathbf{r}} &=&\frac{1}{rR}\left[
\left( \nabla \psi \cdot \widehat{\psi }\right) \cos \alpha \frac{\partial }{
\partial \psi }\left[ rRS_{\mathbf{r}}^{r}\right] +\sin \alpha \frac{%
\partial }{\partial s}\left[ rRS_{\mathbf{r}}^{r}\right] \right]  \nonumber
\\
&&+\frac{1}{r^{2}R}\left[ -\left( \nabla \psi \cdot \widehat{\psi }\right)
r\sin \alpha \frac{\partial }{\partial \psi }\left[ rRS_{\mathbf{r}}^{\theta
}\right] +r\cos \alpha \frac{\partial }{\partial s}\left[ rRS_{\mathbf{r}%
}^{\theta }\right] \right]
\end{eqnarray}

The value of the flux $\psi $ in the center can be arbitrarily set to $0$,
and the evolution of $\psi $ must be monotonic (nested flux-surfaces), so
that the factor $\nabla \psi \cdot \widehat{\psi }$ can be expressed as 
\begin{equation}
\nabla \psi \cdot \widehat{\psi }=\sigma _{\psi }\left| \nabla \psi \right|
\end{equation}
where $\sigma _{\psi }=\pm 1$ is constant across the plasma.

Finally, 
\begin{eqnarray}
\nabla _{\mathbf{r}}\cdot \mathbf{S}_{\mathbf{r}} &=&\frac{1}{rR}\sigma
_{\psi }\left| \nabla \psi \right| \cos \alpha \frac{\partial }{\partial %
\psi }\left[ rR\left( \cos \alpha S_{\mathbf{r}}^{\psi }+\sin \alpha S_{%
\mathbf{r}}^{s}\right) \right]   \nonumber \\
&&+\frac{1}{rR}\sin \alpha \frac{\partial }{\partial s}\left[ rR\left( \cos
\alpha S_{\mathbf{r}}^{\psi }+\sin \alpha S_{\mathbf{r}}^{s}\right) \right] 
\nonumber \\
&&-\frac{1}{rR}\sigma _{\psi }\left| \nabla \psi \right| \sin \alpha \frac{%
\partial }{\partial \psi }\left[ rR\left( -\sin \alpha S_{\mathbf{r}}^{\psi
}+\cos \alpha S_{\mathbf{r}}^{s}\right) \right]   \nonumber \\
&&+\frac{1}{rR}\cos \alpha \frac{\partial }{\partial s}\left[ rR\left( -\sin
\alpha S_{\mathbf{r}}^{\psi }+\cos \alpha S_{\mathbf{r}}^{s}\right) \right] 
\end{eqnarray}
and regrouping, 
\begin{eqnarray}
rR\nabla _{\mathbf{r}}\cdot \mathbf{S}_{\mathbf{r}} &=&\sigma _{\psi }\left|
\nabla \psi \right| \cos \alpha \left( \frac{\partial }{\partial \psi }%
\left[ \cos \alpha rRS_{\mathbf{r}}^{\psi }\right] +\frac{\partial }{%
\partial \psi }\left[ \sin \alpha rRS_{\mathbf{r}}^{s}\right] \right)  
\nonumber \\
&&+\sin \alpha \left( \frac{\partial }{\partial s}\left[ \cos \alpha rRS_{%
\mathbf{r}}^{\psi }\right] +\frac{\partial }{\partial s}\left[ \sin \alpha
rRS_{\mathbf{r}}^{s}\right] \right)   \nonumber \\
&&-\sigma _{\psi }\left| \nabla \psi \right| \sin \alpha \left( -\frac{%
\partial }{\partial \psi }\left[ \sin \alpha rRS_{\mathbf{r}}^{\psi }\right]
+\frac{\partial }{\partial \psi }\left[ \cos \alpha rRS_{\mathbf{r}%
}^{s}\right] \right)   \nonumber \\
&&+\cos \alpha \left( -\frac{\partial }{\partial s}\left[ \sin \alpha rRS_{%
\mathbf{r}}^{\psi }\right] +\frac{\partial }{\partial s}\left[ \cos \alpha
rRS_{\mathbf{r}}^{s}\right] \right) 
\end{eqnarray}
\begin{eqnarray}
rR\nabla _{\mathbf{r}}\cdot \mathbf{S}_{\mathbf{r}} &=&\sigma _{\psi }\left|
\nabla \psi \right| \cos \alpha \left( \cos \alpha \frac{\partial }{\partial %
\psi }\left[ rRS_{\mathbf{r}}^{\psi }\right] +rRS_{\mathbf{r}}^{\psi }\frac{%
\partial }{\partial \psi }\left[ \cos \alpha \right] \right)   \nonumber \\
&&+\sigma _{\psi }\left| \nabla \psi \right| \cos \alpha \left( \sin \alpha 
\frac{\partial }{\partial \psi }\left[ rRS_{\mathbf{r}}^{s}\right] +rRS_{%
\mathbf{r}}^{s}\frac{\partial }{\partial \psi }\left[ \sin \alpha \right]
\right)   \nonumber \\
&&+\sigma _{\psi }\left| \nabla \psi \right| \sin \alpha \left( \sin \alpha 
\frac{\partial }{\partial \psi }\left[ rRS_{\mathbf{r}}^{\psi }\right] +rRS_{%
\mathbf{r}}^{\psi }\frac{\partial }{\partial \psi }\left[ \sin \alpha
\right] \right)   \nonumber \\
&&-\sigma _{\psi }\left| \nabla \psi \right| \sin \alpha \left( \cos \alpha 
\frac{\partial }{\partial \psi }\left[ rRS_{\mathbf{r}}^{s}\right] +rRS_{%
\mathbf{r}}^{s}\frac{\partial }{\partial \psi }\left[ \cos \alpha \right]
\right)   \nonumber \\
&&+\sin \alpha \left( \cos \alpha \frac{\partial }{\partial s}\left[ rRS_{%
\mathbf{r}}^{\psi }\right] +rRS_{\mathbf{r}}^{\psi }\frac{\partial }{%
\partial s}\left[ \cos \alpha \right] \right)   \nonumber \\
&&+\sin \alpha \left( \sin \alpha \frac{\partial }{\partial s}\left[ rRS_{%
\mathbf{r}}^{s}\right] +rRS_{\mathbf{r}}^{s}\frac{\partial }{\partial s}%
\left[ \sin \alpha \right] \right)   \nonumber \\
&&-\cos \alpha \left( \sin \alpha \frac{\partial }{\partial s}\left[ rRS_{%
\mathbf{r}}^{\psi }\right] +rRS_{\mathbf{r}}^{\psi }\frac{\partial }{%
\partial s}\left[ \sin \alpha \right] \right)   \nonumber \\
&&+\cos \alpha \left( \cos \alpha \frac{\partial }{\partial s}\left[ rRS_{%
\mathbf{r}}^{s}\right] +rRS_{\mathbf{r}}^{s}\frac{\partial }{\partial s}%
\left[ \cos \alpha \right] \right) 
\end{eqnarray}
\begin{eqnarray}
rR\nabla _{\mathbf{r}}\cdot \mathbf{S}_{\mathbf{r}} &=&\sigma _{\psi }\left|
\nabla \psi \right| \frac{\partial }{\partial \psi }\left[ rRS_{\mathbf{r}%
}^{\psi }\right]   \nonumber \\
&&+\frac{\partial }{\partial s}\left[ rRS_{\mathbf{r}}^{s}\right]   \nonumber
\\
&&+\sigma _{\psi }\left| \nabla \psi \right| rRS_{\mathbf{r}}^{\psi }\left[
\cos \alpha \frac{\partial \cos \alpha }{\partial \psi }+\sin \alpha \frac{%
\partial \sin \alpha }{\partial \psi }\right]   \nonumber \\
&&+\sigma _{\psi }\left| \nabla \psi \right| rRS_{\mathbf{r}}^{s}\left[ \cos
\alpha \frac{\partial \sin \alpha }{\partial \psi }-\sin \alpha \frac{%
\partial \cos \alpha }{\partial \psi }\right]   \nonumber \\
&&+rRS_{\mathbf{r}}^{\psi }\left[ \sin \alpha \frac{\partial \cos \alpha }{%
\partial s}-\cos \alpha \frac{\partial \sin \alpha }{\partial s}\right]  
\nonumber \\
&&+rRS_{\mathbf{r}}^{s}\left[ \sin \alpha \frac{\partial \sin \alpha }{%
\partial s}+\cos \alpha \frac{\partial \cos \alpha }{\partial s}\right] 
\end{eqnarray}
\begin{eqnarray}
rR\nabla _{\mathbf{r}}\cdot \mathbf{S}_{\mathbf{r}} &=&\sigma _{\psi }\left|
\nabla \psi \right| \frac{\partial }{\partial \psi }\left[ rRS_{\mathbf{r}%
}^{\psi }\right] +\frac{\partial }{\partial s}\left[ rRS_{\mathbf{r}%
}^{s}\right]   \nonumber \\
&&+\sigma _{\psi }\left| \nabla \psi \right| rRS_{\mathbf{r}}^{s}\frac{%
\partial \alpha }{\partial \psi }-rRS_{\mathbf{r}}^{\psi }\frac{\partial %
\alpha }{\partial s}
\end{eqnarray}
\begin{eqnarray}
\nabla _{\mathbf{r}}\cdot \mathbf{S}_{\mathbf{r}} &=&\sigma _{\psi }\frac{%
\left| \nabla \psi \right| }{rR}\frac{\partial }{\partial \psi }\left[ rRS_{%
\mathbf{r}}^{\psi }\right] +\frac{1}{rR}\frac{\partial }{\partial s}\left[
rRS_{\mathbf{r}}^{s}\right]   \nonumber \\
&&+\sigma _{\psi }\left| \nabla \psi \right| S_{\mathbf{r}}^{s}\frac{%
\partial \alpha }{\partial \psi }-S_{\mathbf{r}}^{\psi }\frac{\partial %
\alpha }{\partial s}
\end{eqnarray}

The bounce average operation is 
\begin{eqnarray}
\left\{ \mathcal{A}\right\}  &=&\frac{1}{\tau _{b}}\left[ \frac{1}{2}%
\sum_{\sigma }\right] _{T}\int_{s_{\min }}^{s_{\max }}\frac{ds}{\left|
v_{s}\right| }\mathcal{A}  \nonumber \\
&=&\frac{1}{\tau _{b}}\left[ \frac{1}{2}\sum_{\sigma }\right]
_{T}\int_{s_{\min }}^{s_{\max }}\frac{ds}{\left| v_{\Vert }\right| }\frac{B}{%
B_{P}}\mathcal{A}
\end{eqnarray}

Because the parallel dynamics is described in our model, we know the fluxes
along the magnetic field lines, which is simply
\begin{equation}
S_{\mathbf{r}}^{s}=v_{s}f\left( v_{s}\right) 
\end{equation}

Since $f$ is constant along a field line, $\partial f/\partial s=0$ and we
have simply
\begin{eqnarray}
\left\{ \frac{1}{rR}\frac{\partial }{\partial s}\left[ rRS_{\mathbf{r}%
}^{s}\right] \right\}  &=&\frac{1}{\tau _{b}}\left[ \frac{1}{2}\sum_{\sigma
}\right] _{T}\int_{s_{\min }}^{s_{\max }}\frac{ds}{\left| v_{s}\right| }%
\frac{1}{rR}\frac{\partial }{\partial s}\left[ rRS_{\mathbf{r}}^{s}\right]  
\nonumber \\
&=&f\frac{1}{\tau _{b}}\left[ \frac{1}{2}\sum_{\sigma }\right] _{T}\sigma
\int_{s_{\min }}^{s_{\max }}\frac{ds}{v_{s}}\frac{1}{rR}\frac{\partial }{%
\partial s}\left[ rRv_{s}\right]  \\
&=&f\frac{1}{\tau _{b}}\left[ \frac{1}{2}\sum_{\sigma }\right] _{T}\sigma
\int_{s_{\min }}^{s_{\max }}ds\frac{\partial }{\partial s}\ln \left[ \frac{%
rRv_{s}}{r\left( \psi ,0\right) R\left( \psi ,0\right) v_{s}\left( \psi
,0,v_{s0}\right) }\right]  \\
&=&f\frac{1}{\tau _{b}}\left[ \frac{1}{2}\sum_{\sigma }\right] _{T}\sigma
\left( \ln \left[ \frac{r\left( \psi ,s_{\max }\right) }{r\left( \psi
,s_{\min }\right) }\right] \ln \left[ \frac{R\left( \psi ,s_{\max }\right) }{%
R\left( \psi ,s_{\min }\right) }\right] +\ln \left[ \frac{v_{s}\left( \psi
,s_{\max },v_{s0}\right) }{v_{s}\left( \psi ,s_{\min },v_{s0}\right) }%
\right] \right) 
\end{eqnarray}

For passing particles, the positions $\left( \psi ,s_{\min }\right) $ and $%
\left( \psi ,s_{\max }\right) $ coincide and therefore the term vanishes.
Because the factor is an odd function of $\sigma $, the sum over $\sigma $
for trapped particle leads to $0$. Consequently,
\begin{equation}
\left\{ \frac{1}{rR}\frac{\partial }{\partial s}\left[ rRS_{\mathbf{r}%
}^{s}\right] \right\} =0
\end{equation}

Using 
\begin{equation}
\left| \nabla \psi \right| =RB_{P}
\end{equation}
we find 
\begin{equation}
\left\{ \sigma _{\psi }\frac{\left| \nabla \psi \right| }{rR}\frac{\partial 
}{\partial \psi }\left[ rRS_{\mathbf{r}}^{\psi }\right] \right\} =\sigma
_{\psi }\frac{1}{\tau _{b}}\left[ \frac{1}{2}\sum_{\sigma }\right]
_{T}\int_{s_{\min }}^{s_{\max }}\frac{ds}{\left| v_{\Vert }\right| }\frac{B}{%
r}\frac{\partial }{\partial \psi }\left[ rRS_{\mathbf{r}}^{\psi }\right] 
\end{equation}
\begin{equation}
\left\{ \sigma _{\psi }\left| \nabla \psi \right| S_{\mathbf{r}}^{s}\frac{%
\partial \alpha }{\partial \psi }\right\} =\sigma _{\psi }\frac{1}{\tau _{b}}%
\left[ \frac{1}{2}\sum_{\sigma }\right] _{T}\int_{s_{\min }}^{s_{\max }}%
\frac{ds}{\left| v_{\Vert }\right| }RBS_{\mathbf{r}}^{s}\frac{\partial %
\alpha }{\partial \psi }
\end{equation}
\begin{equation}
\left\{ -S_{\mathbf{r}}^{\psi }\frac{\partial \alpha }{\partial s}\right\} =-%
\frac{1}{\tau _{b}}\left[ \frac{1}{2}\sum_{\sigma }\right] _{T}\int_{s_{\min
}}^{s_{\max }}\frac{ds}{\left| v_{s}\right| }S_{\mathbf{r}}^{\psi }\frac{%
\partial \alpha }{\partial s}
\end{equation}

Cylindrincal limit with circular magnetic flux surfaces

In the cylindrical limit with circular concentric magnetic flux surfaces $%
R=R_{0},$ and $\alpha =0.$ Furthermore, since 
\begin{equation}
\left. \frac{\partial }{\partial r}\right| _{s}=\frac{\nabla \psi \cdot 
\widehat{\psi }}{\nabla r\cdot \widehat{r}}\left( \widehat{\psi }\cdot 
\widehat{r}\right) \frac{\partial }{\partial \psi }=\left( \nabla \psi \cdot 
\widehat{\psi }\right) \frac{\partial }{\partial \psi }=\sigma _{\psi
}\left| \nabla \psi \right| \frac{\partial }{\partial \psi }
\end{equation}
one obtains 
\begin{equation}
\left\{ \sigma _{\psi }\frac{\left| \nabla \psi \right| }{rR}\frac{\partial 
}{\partial \psi }\left[ rRS_{\psi }\right] \right\} =\frac{1}{r}\frac{%
\partial }{\partial r}\left[ r\left\{ S_{\mathbf{r}}^{\psi }\right\} \right]
=\frac{1}{r}\frac{\partial }{\partial r}\left[ r\left\{ S_{\mathbf{r}%
}^{r}\right\} \right]
\end{equation}
\[
\left\{ \frac{1}{rR}\frac{\partial }{\partial s}\left[ rRS_{\mathbf{r}%
}^{s}\right] \right\} =0 
\]
\begin{equation}
\left\{ \sigma _{\psi }\left| \nabla \psi \right| S_{\mathbf{r}}^{s}\frac{
\partial \alpha }{\partial \psi }\right\} =0
\end{equation}
\begin{equation}
\left\{ -S_{\mathbf{r}}^{\psi }\frac{\partial \alpha }{\partial s}\right\} =0
\end{equation}
The poloidal magnetic field $B_{p}$ is then simply, 
\begin{equation}
B_{p}=\frac{1}{R_{0}}\frac{d\psi \left( r\right) }{dr}=\frac{\psi ^{\prime }%
}{R_{0}}
\end{equation}
and the expression \ref{qtilde} is the well known local value of the
cylindrical safety factor

\begin{equation}
\widetilde{q}\left( \psi \right) =\frac{r}{R_{0}}\frac{B_{0}}{B_{p}}=r\frac{
B_{0}}{\psi ^{\prime }}
\end{equation}

Furthermore, in this simplified configuration, the fraction of trapped
electrons vanishes, and consequently $\theta _{\min }=0,$ $\theta _{\max
}=2\pi ,$ $\frac{\xi _{0}}{\xi }=1$ and $\lambda =1.$ Hence, using the fact
that $S_{\psi }^{\psi }=S_{\psi }^{r}.$
